\section{Projekt systemu}

W niniejszym rozdziale przedstawiono projekt systemu informatycznego służącego do analizy techniki martwego ciągu na podstawie nagrania wideo. Rozdział stanowi przejście pomiędzy analizą wymagań a opisem implementacji. Z tego względu skupiono się w nim na ogólnej organizacji aplikacji, podziale odpowiedzialności pomiędzy jej główne części oraz sposobie przepływu danych od przesłania nagrania do wyświetlenia wyniku.

System został zaprojektowany jako aplikacja sieciowa składająca się z trzech głównych części: aplikacji klienckiej, aplikacji serwerowej oraz modułu analitycznego zaimplementowanego w języku Python. W dalszej części pracy aplikacja kliencka określana jest również jako frontend, natomiast aplikacja serwerowa jako backend. Pojęcia te stosowane są ze względu na ich powszechne wykorzystanie w projektach aplikacji sieciowych oraz bezpośrednie odniesienie do struktury zaimplementowanego projektu.

Frontend odpowiada za kontakt z użytkownikiem, czyli wybór pliku wideo, przesłanie go do analizy oraz prezentację wyniku. Backend przyjmuje nagranie, zapisuje informacje o analizie i uruchamia moduł analityczny. Moduł analityczny przetwarza nagranie oraz zwraca wynik w postaci danych, które mogą zostać zapisane po stronie serwera i przedstawione w interfejsie użytkownika.

\subsection{Ogólna struktura systemu}

Podstawowym założeniem projektowym było rozdzielenie obsługi użytkownika od analizy nagrania. Użytkownik nie uruchamia bezpośrednio programu analizującego wideo. Korzysta jedynie z aplikacji sieciowej, za pośrednictwem której wybiera plik i otrzymuje wynik. Backend pełni w tym układzie rolę pośrednika pomiędzy interfejsem użytkownika a modułem analitycznym.

Na poziomie projektu wyróżniono następujące części systemu:

\begin{itemize}
\item \textbf{frontend} -- aplikacja kliencka uruchamiana w przeglądarce internetowej,
\item \textbf{backend} -- aplikacja serwerowa odpowiedzialna za obsługę żądań i zapisywanie danych analizy,
\item \textbf{moduł analityczny} -- wydzielona część systemu odpowiedzialna za właściwą analizę nagrania.
\end{itemize}

Taki podział odpowiada rzeczywistej strukturze projektu. Frontend i backend są oddzielnymi częściami aplikacji sieciowej, natomiast moduł analityczny znajduje się poza kodem odpowiedzialnym za obsługę żądań HTTP. Dzięki temu kod związany z analizą ruchu jest oddzielony od kodu odpowiedzialnego za interfejs użytkownika oraz komunikację sieciową.

\newpage
Ogólną architekturę systemu przedstawiono na rysunku~\ref{fig:diagram-architektury-systemu}.

\begin{figure}[H]
\centering
\includegraphics[width=0.95\textwidth]{figures/diagram_architektury_systemu.png}
\caption{Architektura systemu analizy techniki martwego ciągu. \newline Źródło: opracowanie własne.}
\label{fig:diagram-architektury-systemu}
\end{figure}

\subsection{Aplikacja kliencka}

Aplikacja kliencka jest częścią systemu widoczną dla użytkownika. Jej zadaniem jest umożliwienie przesłania nagrania wideo oraz przedstawienie wyniku analizy w czytelnej formie. Użytkownik powinien móc wykonać całą operację bez znajomości szczegółów działania modelu estymacji pozy człowieka oraz bez bezpośredniego kontaktu z modułem analitycznym.

W części klienckiej przewidziano następujące zadania:

\begin{itemize}
\item wyświetlenie krótkiej informacji o sposobie przygotowania nagrania,
\item umożliwienie wyboru pliku wideo,
\item sprawdzenie podstawowych warunków dotyczących wybranego pliku,
\item przesłanie nagrania do backendu,
\item pokazanie informacji o trwającej analizie,
\item wyświetlenie wyniku zwróconego przez serwer,
\item pokazanie komunikatu błędu, jeżeli analiza nie zakończy się poprawnie.
\end{itemize}

Interfejs użytkownika powinien w pierwszej kolejności prezentować wynik ogólny, a dopiero później szczegóły dotyczące poszczególnych powtórzeń. Jest to istotne, ponieważ użytkownik nie powinien być od razu obciążany dużą liczbą szczegółowych danych liczbowych. Szczegółowe informacje mają pełnić funkcję pomocniczą i umożliwiać dokładniejsze przejrzenie wyniku wtedy, gdy użytkownik chce sprawdzić, z czego wynika dana ocena.

\subsection{Aplikacja serwerowa}

Backend odpowiada za przyjęcie nagrania przesłanego z frontendu oraz uruchomienie analizy. Po otrzymaniu żądania tworzy rekord analizy na podstawie przesłanych danych i zapisuje w nim plik wideo. Początkowo analiza otrzymuje status informujący o oczekiwaniu na rozpoczęcie przetwarzania. Następnie status zostaje zmieniony na informujący o rozpoczęciu analizy, a backend przekazuje ścieżkę zapisanego pliku do modułu analitycznego.

Dla każdej przesłanej analizy przechowywany jest status informujący, czy analiza oczekuje na wykonanie, jest przetwarzana, zakończyła się poprawnie czy zakończyła się błędem. Oprócz statusu przechowywany jest także plik wideo, wynik analizy, ewentualny komunikat błędu oraz czas rozpoczęcia i zakończenia przetwarzania.

Do głównych zadań backendu należą:

\begin{itemize}
\item odebranie pliku przesłanego przez użytkownika,
\item zapisanie nagrania po stronie serwera,
\item utworzenie rekordu analizy,
\item ustawienie odpowiedniego statusu analizy,
\item uruchomienie modułu analitycznego,
\item zapisanie wyniku zwróconego przez moduł analityczny,
\item zwrócenie odpowiedzi do aplikacji klienckiej.
\end{itemize}

Backend nie wykonuje samodzielnie obliczeń biomechanicznych. Nie wykrywa punktów kluczowych, nie dzieli nagrania na powtórzenia i nie ocenia techniki. Zadania te zostały przypisane modułowi analitycznemu. Dzięki temu część serwerowa pozostaje odpowiedzialna głównie za obsługę żądań, zapis danych oraz komunikację z modułem analitycznym.

\subsection{Moduł analityczny DO OCENY}

Moduł analityczny jest częścią systemu odpowiedzialną za przetworzenie nagrania wideo. Z punktu widzenia implementacji jest to wydzielona część projektu napisana w języku Python, uruchamiana z poziomu backendu po przesłaniu nagrania przez użytkownika. Wejściem modułu jest ścieżka do pliku wideo, a jego wyjściem jest struktura danych zawierająca wynik analizy.

W module analitycznym wykorzystano gotowy model estymacji pozy człowieka \cite{mediapipe_pose_landmarker}. Model ten odpowiada za wykrywanie punktów kluczowych sylwetki w kolejnych klatkach nagrania. Dalsze przetwarzanie realizowane jest przez kod przygotowany w ramach projektu. Obejmuje ono wybór potrzebnych punktów, obliczenie parametrów opisujących ruch, wykrycie powtórzeń oraz ocenę techniki.

W module analizy wideo zastosowano wariant Heavy modelu MediaPipe Pose Landmarker. Spośród dostępnych wariantów Lite, Full i Heavy charakteryzuje się on największym kosztem obliczeniowym, ale jednocześnie osiąga najlepsze wyniki estymacji punktów kluczowych w testach przedstawionych przez twórców rozwiązania \cite{mediapipe_pose_docs}. Wybór tego wariantu wynika z charakteru projektowanego systemu. Analiza nie jest wykonywana w czasie rzeczywistym. Użytkownik najpierw przesyła nagranie do aplikacji, a następnie materiał jest przetwarzany po stronie serwera. Czas przetwarzania pojedynczej klatki nie stanowi więc najważniejszego kryterium wyboru modelu. Większe znaczenie ma uzyskanie możliwie dobrej jakości estymacji punktów, ponieważ ich współrzędne stanowią dane wejściowe dla kolejnych etapów analizy biomechanicznej.

Przyjęto również, że wybrane klatki nagrania są przekazywane do modelu jako niezależne obrazy. MediaPipe Pose Landmarker umożliwia pracę w trybie przeznaczonym dla materiału wideo, w którym informacje z poprzednich klatek mogą być wykorzystywane do śledzenia sylwetki w kolejnych klatkach \cite{mediapipe_pose_landmarker,bazarevsky2020blazepose}. W opracowanym systemie zastosowano natomiast tryb pojedynczego obrazu, dzięki czemu wynik estymacji dla analizowanej klatki nie zależy od poprawności śledzenia sylwetki w klatkach wcześniejszych.

Decyzja ta wiąże się z większym kosztem obliczeniowym, ponieważ model nie wykorzystuje mechanizmu śledzenia do ograniczenia liczby ponownych detekcji. Koszt ten został zaakceptowany ze względu na sposób działania aplikacji oraz brak wymagania pracy w czasie rzeczywistym. Jednocześnie nie przyjęto założenia, że tryb pojedynczego obrazu jest sam w sobie dokładniejszy od trybu wideo. Celem tej decyzji było niezależne przetwarzanie kolejnych próbek obrazu i uniknięcie uzależnienia wyniku bieżącej klatki od stanu śledzenia wyznaczonego na podstawie klatek wcześniejszych.

Proces działania modułu analitycznego można opisać następująco:

\begin{itemize}
\item wczytanie nagrania wideo,
\item przetwarzanie kolejnych klatek nagrania,
\item wykrycie punktów kluczowych sylwetki,
\item zapisanie danych o położeniu punktów w czasie,
\item wyznaczenie pozycji sztangi na podstawie wybranych punktów,
\item obliczenie wybranych parametrów biomechanicznych,
\item wykrycie powtórzeń w nagraniu,
\item przygotowanie podsumowania każdego powtórzenia,
\item wykonanie oceny techniki,
\item zwrócenie wyniku do backendu.
\end{itemize}

Wewnątrz modułu analitycznego można wyróżnić kilka etapów przetwarzania, jednak nie są one osobnymi modułami systemu widocznymi z zewnątrz. Stanowią kolejne kroki wykonywane w ramach jednej części projektu. Szczegółowy opis ich implementacji zostanie przedstawiony w kolejnym rozdziale.

\subsection{Przepływ danych w systemie}

Przepływ danych rozpoczyna się w momencie, gdy użytkownik wybiera plik wideo w aplikacji klienckiej. Plik zostaje przesłany do backendu, gdzie tworzony jest rekord analizy. Backend zapisuje nagranie i przekazuje ścieżkę do pliku do modułu analitycznego. Po zakończeniu analizy wynik zostaje zapisany po stronie serwera i zwrócony do frontendu.

Przepływ danych można przedstawić w następujących krokach:

\begin{enumerate}
\item użytkownik wybiera plik wideo w interfejsie aplikacji,
\item frontend przesyła plik do backendu,
\item backend tworzy rekord analizy i zapisuje przesłane nagranie,
\item backend uruchamia moduł analityczny, przekazując mu ścieżkę do pliku,
\item moduł analityczny przetwarza nagranie i tworzy wynik analizy,
\item backend zapisuje wynik w rekordzie analizy,
\item backend zwraca wynik do aplikacji klienckiej,
\item frontend prezentuje użytkownikowi podsumowanie oraz szczegóły analizy.
\end{enumerate}

W obecnej wersji projektu analiza wykonywana jest bezpośrednio po przesłaniu nagrania. Oznacza to, że użytkownik czeka na odpowiedź zawierającą wynik lub informację o błędzie. Takie rozwiązanie jest wystarczające dla aplikacji przygotowywanej w ramach pracy, ponieważ system analizuje pojedyncze nagranie, a nie obsługuje wielu równoległych użytkowników ani kolejki zadań.

\subsection{Projekt wyniku analizy}

Wynik analizy powinien być możliwy do zapisania po stronie backendu oraz do bezpośredniego wykorzystania przez frontend. Z tego względu przyjęto strukturę danych możliwą do zapisania w formacie JSON. Taki format pozwala przechowywać zarówno proste informacje, takie jak liczba wykrytych powtórzeń, jak i bardziej złożone dane opisujące poszczególne etapy analizy.

Wynik analizy powinien zawierać przede wszystkim:

\begin{itemize}
\item liczbę wykrytych powtórzeń,
\item dane czasowe dotyczące przebiegu nagrania,
\item informacje o wykrytych powtórzeniach,
\item podsumowanie każdego powtórzenia,
\item ocenę techniki dla poszczególnych powtórzeń,
\item ocenę spójności wykonania całej serii.
\end{itemize}

Każde powtórzenie traktowane jest jako osobny element wyniku. Dzięki temu aplikacja kliencka może pokazać listę powtórzeń i przy każdym z nich przedstawić status oraz komunikaty dotyczące techniki. Takie rozwiązanie dobrze pasuje do charakteru ćwiczenia, ponieważ jedna seria może zawierać zarówno powtórzenia poprawne, jak i takie, które wymagają uwagi.

Wynik nie powinien ograniczać się tylko do jednej końcowej oceny. Bardziej użyteczne jest pokazanie użytkownikowi, które powtórzenia zostały wykonane poprawnie, a które odbiegały od przyjętych kryteriów. Dzięki temu użytkownik może zauważyć, czy problem dotyczy całej serii, czy tylko pojedynczych powtórzeń.

\subsection{Obsługa błędów}

W projekcie uwzględniono sytuacje, w których analiza nie może zostać wykonana poprawnie. Błąd może pojawić się zarówno po stronie aplikacji webowej, jak i w czasie działania modułu analitycznego. Z tego względu backend powinien przechowywać nie tylko wynik, ale także status analizy oraz komunikat błędu.

Do przykładowych sytuacji błędnych należą:

\begin{itemize}
\item brak wybranego pliku po stronie użytkownika,
\item niepoprawny lub nieobsługiwany plik wideo,
\item problem z zapisaniem przesłanego nagrania,
\item brak możliwości odczytania klatek nagrania,
\item brak wykrycia sylwetki w nagraniu,
\item brak wykrytych powtórzeń,
\item błąd podczas działania modułu analitycznego.
\end{itemize}

W przypadku błędu system powinien zakończyć analizę w sposób kontrolowany. Backend zapisuje wtedy informację o niepowodzeniu oraz komunikat błędu. Frontend powinien zaprezentować ten komunikat użytkownikowi w prosty sposób, bez pokazywania szczegółów technicznych, które nie są potrzebne do zrozumienia problemu.

\subsection{Uzasadnienie przyjętego rozwiązania}

Przyjęty podział na frontend, backend i moduł analityczny jest zgodny z charakterem systemu. Frontend pozostaje odpowiedzialny za interfejs użytkownika, backend za obsługę żądań i zapis danych, a moduł analityczny za przetwarzanie nagrania. Dzięki temu każda część systemu ma jasno określoną odpowiedzialność.

Wydzielenie modułu analitycznego w osobnym katalogu projektu ułatwia jego dalszy rozwój. Możliwe jest dodawanie nowych parametrów biomechanicznych, zmiana sposobu wykrywania powtórzeń lub modyfikacja reguł oceny techniki bez konieczności przebudowy interfejsu użytkownika. Jednocześnie backend może korzystać z tego modułu przez jedno miejsce odpowiedzialne za uruchomienie analizy i przygotowanie wyniku do zapisu.

Tak zaprojektowana architektura spełnia wymagania określone w poprzednim rozdziale. System umożliwia przesłanie nagrania, wykonanie analizy po stronie serwera oraz prezentację wyników użytkownikowi. Jednocześnie nie sugeruje istnienia dodatkowych niezależnych usług, których w projekcie faktycznie nie ma.
